
This paper reports two diagnostic experiments investigating whether max-pooling aggregation over query-head scores within GQA groups (GQA-GroupMax) produces meaningfully different KV cache eviction decisions than mean-pooling (GQA-GroupMean). The existence experiment (H-E1) found that within-group query-head attention cosine similarity falls below 0.9 in 40.6\% of LLaMA-3-8B layers, below the 50\% pre-specified threshold for the primary model, though Mistral-7B (62.5\%) and Phi-3-mini (78.1\%) exceed the threshold. The mechanism experiment (H-M1) found a global median Spearman rank correlation of 0.962 between GroupMax and GroupMean eviction scores on LLaMA-3-8B at 2048-token context length, above the 0.95 threshold for meaningful divergence.

Both gate failures are quantitative near-misses associated with an identified confound: the sequence length truncation required by memory constraints reduces the attention sparsity that motivates the GroupMax/GroupMean distinction. The mechanism is confirmed active (non-zero divergence, per-layer hook verification), and 7 of 32 layers show per-layer correlation below the 0.95 threshold.

The formalization of the aggregation function as an explicit design choice for GQA-native KV cache eviction remains a contribution independent of these gate outcomes: no prior eviction paper identifies or studies this choice. The infrastructure built and validated here provides a ready foundation for downstream performance experiments (H-M2, H-M3) with full-context sequence lengths, which are the necessary next step to evaluate whether the GroupMax/GroupMean distinction has practical significance for LongBench retrieval task performance.

